\section{Knowledge Graph Construction}
\label{sec:construction}

The LACK knowledge graph is produced by a four-phase pipeline that transforms the relation extraction outputs described in Section~\ref{sec:extraction} into a fully materialised, ontology-aligned RDF graph.
Each phase is implemented as a SPARQL or SPARQL-Anything~\cite{sparqlanything} query, executed with the \texttt{fx} command-line tool, making the entire pipeline reproducible from the raw CSV inputs.

\subsection{Phase 1: Relation Alignment}
\label{sec:construction:alignment}

The relation extraction pipeline produces relations expressed as natural-language surface forms (e.g.\ \emph{is a member of}, \emph{has funded}).
The first phase enumerates the distinct relation types appearing across both the Desmog and LobbyMap extraction outputs --- 152 types in total, covering 78,734 relation instances --- and aligns them manually to LACK ontology properties.
The result is a mapping table that associates each surface-form relation type with a canonical ontology IRI (e.g.\ \texttt{lack:memberOf}, \texttt{lack:fundedBy}), which is consumed by all subsequent phases.
This manual alignment step is where domain judgment is applied: surface forms that are ambiguous, too infrequent, or not reliably mappable to any ontology property are flagged and excluded from the graph.

\subsection{Phase 2: Entity URI Generation}
\label{sec:construction:entities}

Stable, dereferenceable IRIs are minted for each entity using a priority scheme over available external identifiers:

\begin{enumerate}
    \item Wikidata QID
    \item Wikipedia page URL
    \item Web page URL
    \item Surface form + evidence link (fallback)
\end{enumerate}

\noindent Each entity IRI takes the form \texttt{lack-entity:\textit{SHA1}}, where the hash is computed from the highest-priority identifier available for that entity.
This scheme ensures that entities with a known Wikidata entry receive a stable, semantically grounded IRI regardless of how they are named in the source data, while still providing a consistent IRI for entities that have no external presence.
Each entity is assigned a top-level type (\texttt{lack:Person} or \texttt{lack:Collective}), a granular subtype (e.g.\ \texttt{lack-type:think\_tank}, \texttt{lack-type:company}), an \texttt{rdfs:label}, and where available an \texttt{owl:sameAs} link to the corresponding Wikidata entity.

\subsection{Phase 3: Relation Graph Construction}
\label{sec:construction:relations}

With entity IRIs in place, the relation triples are constructed from the raw extraction outputs and the relation alignment mappings produced in Phase~1.
A na\"ive join of all inputs would exceed memory constraints, so the phase is split into two steps: first, an entity index is built in memory (mapping surface-form + evidence-link pairs to entity IRIs); second, the index is used to rewrite each extracted relation as a triple using the appropriate LACK ontology property.

Each relation is represented following the pattern established in the ontology (Section~\ref{sec:ontology}): a direct typed triple connects source and target entities, while a reified \texttt{rdf:Statement} carries the evidence URL (\texttt{dct:source}) and, where available, temporal bounds (\texttt{lack:since}, \texttt{lack:until}).
The same logical relation may be attested by multiple evidence sources, yielding multiple reified instances for a single direct triple; the final graph contains 65,527 distinct relation triples and 77,168 reified instances.

\subsection{Phase 4: Provenance}
\label{sec:construction:provenance}

A dedicated provenance query enriches each reified statement with PROV-O~\cite{provo} annotations.
Each statement is linked to a shared \texttt{prov:Activity} representing the LACK extraction activity, and to a \texttt{prov:Agent} identifying the extraction pipeline, via \texttt{prov:wasGeneratedBy}. \texttt{prov:used} arcs collect all source web pages used in any \texttt{rdf:Statement} produced so far.
%This step is kept separate from Phase~3 to maintain a clean separation between the structural construction of the graph and the attachment of process-level metadata.

\subsection{Phase 5: OWL Inference Materialisation}
\label{sec:construction:inference}

Rather than relying on a full OWL reasoner at query time, we materialise the entailments of its ontology as explicit triples using a sequence of SPARQL CONSTRUCT queries.
%This ensures that all queries over the graph are complete with respect to the ontology's axioms without requiring inference-capable infrastructure.
%The materialisation pipeline proceeds in three steps.
First, inverse property pairs are materialised (e.g.\ \texttt{memberOf} / \texttt{hasMember}, \texttt{fundedBy} / \texttt{hasFunder}) for all ten inverse pairs defined in the ontology.
Second, sub-property entailments are propagated: \texttt{leadsAt} triples are promoted to \texttt{employedBy} triples (and correspondingly for their inverses), and all direct sub-properties of \texttt{associatedWith} generate \texttt{associatedWith} triples.
Third, symmetry is closed for \texttt{associatedWith} and \texttt{hasPartner}.
The steps are executed in this order so that triples generated in earlier steps are available as input to later ones, ensuring completeness of the inference closure.
After materialisation, the graph grows from approximately 66,000 relation triples to over 264,000.